% !Mode:: "TeX:UTF-8"

%%% 说明: 此部分需要自己填写的内容:  (1) 中文摘要及关键词 (2) 英文摘要及关键词
%--------------------------------------------------------------------------
%%%页眉设置
%
%\pagestyle{fancyplain}
%\fancyhf{}  %清除以前对页眉页脚的设置
%\renewcommand{\headrulewidth}{0.5pt}
%\fancyfoot[C]{-\,\thepage\,-}
%\fancypagestyle{plain}
%{
%\fancyhead[CE]{武汉大学博士学位论文}
%\fancyhead[CO]{}
%\renewcommand{\headrulewidth}{0.5pt}
%\fancyfoot{}                                    % clear all footer fields
%\fancyfoot[C]{-\,\thepage\,-}
%}

%%%%%%%%%%%%%%%%%%%%%%%
%%% ------------ 中文摘要 ---------------%%%
%%%%%%%%%%%%%%%%%%%%%%%
\begin{cnabstract}

本模板面向中文学位论文写作，提供可复用的章节组织方式与可配置的排版体系，旨在帮助用户聚焦于内容写作而非排版及其他简单重复性工作。项目地址为 \url{https://github.com/yanboc/yanboc-thesis}。
此模板基于黄正华老师上传至 Overleaf 的\href{https://www.overleaf.com/latex/templates/wu-yi-da-xue-bo-shi-lun-wen-latex-mo-ban/rcdzgvqgkddk}{武汉大学学位论文 LaTeX 模板}，且参考了 \href{https://github.com/whutug/whu-thesis}{WHUTUG 模板}，在此表示感谢。

配置层面，模板以 \texttt{settings/\_\_settings\_\_.tex}
作为总入口，通过 \texttt{\textbackslash input} 命令加载参考文献样式、中文交叉引用、字体与配色、术语系统、绘图样式、快捷命令与快速编译等模块。
这种分层组织方式能够降低耦合度，方便用户在不破坏整体结构的前提下完成局部定制。

模板的具体配置文件被统一放在 \texttt{settings/} 目录下，便于用户集中管理、按需加载。
其中，\texttt{bib\_style.tex}
用于设置参考文献样式；
\texttt{cleveref-chinese.tex} 用于定义图、表、公式、章节等对象的中文引用名称与格式；
\texttt{design.tex} 统一字体、颜色等样式；
\texttt{plotting.tex} 统一 TikZ 与 PGFPlots 的绘图库和图元样式；
\texttt{shortcuts.tex} 提供常用数学符号与算子命令，提升写作效率与表达一致性；
\texttt{quick\_compile.tex} 通过占位图机制减少编译开销，适用于高频迭代阶段。

此外，模板支持基于 \texttt{glossaries} 宏包的术语表与缩略语管理。
\texttt{gls.tex} 配合术语定义文件可实现术语首现、索引输出与术语列表统一管理，便于长文档保持术语表达一致。
在参考文献管理方面，我们使用了 \href{https://github.com/zepinglee/gbt7714-bibtex-style}{GB/T 7714—2015宏包}，并结合 \texttt{natbib} 宏包实现了国标参考文献样式与条目间距调节。

总体而言，本模板强调“结构稳定、配置集中、内容可替换”三项原则。
用户可在保留模板骨架的基础上，将个性化内容逐步填入章节与附录，从而在满足学校格式要求的同时提升写作与维护效率。

\end{cnabstract}
\vspace{1em}\par\vfill

%%%--------- 关键词 -------- %%%
\cnkeywords{论文模板；LaTeX}

%%%%%%%%%%%%%%%%%%%%%%%


%%%%%%%%%%%%%%%%%%%%%%%
%%% ------------ 英文摘要 ---------------%%%
%%%%%%%%%%%%%%%%%%%%%%%

\begin{enabstract}
This template is designed for Chinese academic thesis writing, with a reusable document structure and a centralized configuration system.
It is adapted from the Wuhan University thesis template by Prof. Huang Zhenghua (Overleaf) and also references the WHUTUG template design.
Project address: \url{https://github.com/yanboc/yanboc-thesis}.

The configuration entry is \texttt{settings/\_\_settings\_\_.tex}, which loads modular files for bibliography style, Chinese cross-reference formatting, font and color settings, glossary management, plotting, shortcuts, and quick compilation.
Key modules include \texttt{bib\_style.tex}, \texttt{cleveref-chinese.tex}, \texttt{design.tex}, \texttt{plotting.tex}, \texttt{shortcuts.tex}, \texttt{quick\_compile.tex}, and \texttt{gls.tex}.

The template supports GB/T 7714 bibliography formatting (with \texttt{natbib}) and glossary workflows (with \texttt{glossaries}).
Overall, it follows three principles: stable structure, centralized configuration, and replaceable content, so users can focus on writing while keeping format control.
\end{enabstract}
\vspace{1em}\par\vfill

%%%------ 英文关键词 ------- %%%
\enkeywords{Thesis Template, LaTeX, Typesetting Configuration, Bibliography Management, Chinese Cross-Reference, Quick Compilation, Glossary}
